\chapter{Einleitung}
\label{ch:einleitung}

\section{Motivation}
Wissenschaftliches Schreiben lebt von sauberer Typografie. Seit \citet{knuth1984texbook}
das Satzsystem \TeX{} vorgestellt hat und \citet{lamport1994latex} es mit \LaTeX{}
zugänglich machte, ist es der Standard für technische Abschlussarbeiten.

Neue Quellen fügst du im Tab \enquote{Bibliothek} hinzu -- per DOI, arXiv-ID,
Titelsuche oder einfach durch Importieren eines PDFs. Danach genügt
\verb|\cite{|, und die Autovervollständigung schlägt passende Einträge vor.

\section{Forschungsfrage}
\begin{enumerate}[label=\textbf{F\arabic*}]
  \item Wie lässt sich \dots{}?
  \item Welchen Einfluss hat \dots{}?
\end{enumerate}

\section{Aufbau der Arbeit}
\Cref{ch:grundlagen} führt in die Grundlagen ein, \cref{ch:methodik} beschreibt
das Vorgehen, \cref{ch:ergebnisse} präsentiert die Ergebnisse und
\cref{ch:fazit} schließt mit einem Ausblick.
